\documentclass[11pt]{article}

\usepackage[margin=1in]{geometry}
\usepackage{amsmath,amssymb,amsthm}
\usepackage{graphicx}
\usepackage[hidelinks]{hyperref}

\newtheorem{theorem}{Theorem}
\newtheorem{lemma}{Lemma}
\newtheorem{remark}{Remark}

\newcommand{\APEP}{\operatorname{APEP}}
\newcommand{\Ball}{B}
\newcommand{\norm}[1]{\left\lVert #1 \right\rVert}
\newcommand{\R}{\mathbb{R}}

\title{A Counterexample for Zero APEPs in Infinite \texorpdfstring{$\ell_1$}{ell-1}-Sums}
\author{Automated Math Discovery Smoke Test}
\date{June 8, 2026}

\begin{document}
\maketitle

\section*{Source and Open Question}

This packet concerns Question 8.1 from:

\begin{quote}
Ramon J. Aliaga, Luis C. Garcia-Lirola, Juan Guerrero-Viu, Matias Raja,
Abraham Rueda Zoca, \emph{Almost preserved extreme points},
arXiv:2510.21234.
\end{quote}

The source location is page 37 of the paper.

\begin{center}
\includegraphics[width=\textwidth]{figures/open_problem_crop.png}
\end{center}

Question 8.1 asks:

\begin{quote}
Let $\{X_i:i\in I\}$ be an arbitrary infinite family of Banach spaces and
let $X$ be its $\ell_1$-sum. If $0\in \Ball_X$ is APEP in $\Ball_X$, must
there exist $i$ such that $0\in X_i$ is APEP in $\Ball_{X_i}$?
\end{quote}

We give a negative answer.

\section*{Definitions}

Let $C$ be a bounded closed convex subset of a Banach space $E$. A point
$x\in C$ is an almost preserved extreme point, abbreviated APEP, if every
relatively weakly open subset $W$ of $C$ containing $x$ contains a
non-empty slice of $C$. For the unit ball, a slice has the form
\[
  S(\Ball_E,f,\alpha)
  =
  \{y\in \Ball_E: f(y)>\sup f(\Ball_E)-\alpha\},
\]
where $f\in E^*$ and $\alpha>0$.

\section*{Counterexample}

\begin{theorem}
There is a countably infinite family of Banach spaces $(X_n)$ such that
$0$ is an APEP of the unit ball of $X=(\bigoplus_{n=1}^{\infty}X_n)_1$,
but $0$ is not an APEP of the unit ball of any coordinate $X_n$.
\end{theorem}

\begin{proof}
For each $n\ge 1$, let $X_n=H_n=\R^n$ with its Euclidean norm, and put
\[
  X=\left(\bigoplus_{n=1}^{\infty} H_n\right)_1.
\]
Then
\[
  X^*=\left(\bigoplus_{n=1}^{\infty} H_n\right)_\infty,
\]
where we identify each $H_n^*$ with $H_n$ by the inner product.

First, $0$ is not an APEP of $\Ball_{H_n}$ for any fixed $n$. Indeed, in
finite dimension the weak and norm topologies coincide, so
\[
  W_n=\{h\in \Ball_{H_n}:\norm{h}<1/2\}
\]
is a relatively weakly open neighbourhood of $0$ in $\Ball_{H_n}$. No
non-empty slice of $\Ball_{H_n}$ is contained in $W_n$: if
$S(\Ball_{H_n},f,\alpha)$ is non-empty, then either it is all of
$\Ball_{H_n}$ in the degenerate case, or it contains a unit vector at
which $f$ attains its norm. In both cases it contains a point of norm
$1$.

It remains to show that $0$ is an APEP of $\Ball_X$. Let $W$ be a
relatively weakly open neighbourhood of $0$ in $\Ball_X$. After shrinking
$W$, we may assume that
\[
  W=
  \left\{x\in \Ball_X:
  |x_j^*(x)|<\varepsilon,\quad j=1,\dots,m
  \right\}
\]
for some $x_1^*,\dots,x_m^*\in X^*$ and $\varepsilon>0$. Write
\[
  x_j^*=(a_{j,n})_{n\ge 1},\qquad a_{j,n}\in H_n,
\]
and let
\[
  M=\max_{1\le j\le m}\norm{x_j^*}.
\]
If $M=0$, any non-empty slice of $\Ball_X$ is contained in $W$, so assume
$M>0$.

Choose $n>m$. Since $\dim H_n=n$, there is a unit vector $u\in H_n$
orthogonal to the finite set
\[
  a_{1,n},a_{2,n},\dots,a_{m,n}.
\]
Define $\varphi\in X^*$ by
\[
  \varphi((h_k)_{k\ge 1})=\langle h_n,u\rangle.
\]
Choose $\delta>0$ so small that
\[
  M\bigl(\sqrt{2\delta}+\delta\bigr)<\varepsilon.
\]
Consider the slice
\[
  S=\{x=(h_k)\in\Ball_X:\varphi(x)>1-\delta\}.
\]
This slice is non-empty, since it contains the vector supported only at
coordinate $n$ with value $u$.

We claim that $S\subseteq W$. Let $x=(h_k)\in S$. Since
\[
  1-\delta < \langle h_n,u\rangle\le \norm{h_n}\le
  \sum_{k=1}^{\infty}\norm{h_k}\le 1,
\]
we have
\[
  \sum_{k\ne n}\norm{h_k}<\delta.
\]
Also, using $\norm{h_n}\le 1$ and $\langle h_n,u\rangle>1-\delta$,
\[
  \norm{h_n-u}^2
  =
  \norm{h_n}^2+1-2\langle h_n,u\rangle
  \le 2\delta.
\]
For each $j=1,\dots,m$, orthogonality gives
\[
  \langle a_{j,n},h_n\rangle
  =
  \langle a_{j,n},h_n-u\rangle.
\]
Therefore
\[
\begin{aligned}
  |x_j^*(x)|
  &\le |\langle a_{j,n},h_n\rangle|
      + \sum_{k\ne n}|\langle a_{j,k},h_k\rangle|  \\
  &\le M\norm{h_n-u}+M\sum_{k\ne n}\norm{h_k} \\
  &< M\bigl(\sqrt{2\delta}+\delta\bigr)
   < \varepsilon.
\end{aligned}
\]
Thus $x\in W$ for every $x\in S$, so $S\subseteq W$.

Every relatively weakly open neighbourhood of $0$ in $\Ball_X$ contains
a non-empty slice of $\Ball_X$. Hence $0$ is an APEP of $\Ball_X$, even
though no coordinate has $0$ as an APEP. This gives a negative answer to
Question 8.1.
\end{proof}

\section*{Verification Notes}

The argument is purely analytic. The only quantitative estimate is the
Hilbert-space fact that if $\langle h,u\rangle>1-\delta$ with
$\norm{h}\le 1=\norm{u}$, then $\norm{h-u}\le \sqrt{2\delta}$. The script
\texttt{code/check\_slice\_estimate.py} performs randomized finite
dimensional checks of the orthogonal-vector construction and this bound.
It is a sanity check only; the proof above is the certificate.

\section*{Limitations and Novelty Check}

This packet answers the source question as stated and in its intended
infinite-sum setting: the counterexample is a genuine infinite
$\ell_1$-sum, not a degenerate one-dimensional edge case. I have not
performed a literature search beyond the source paper and the local
pipeline corpus, so the novelty status is ``candidate'' pending human
and literature review.

\end{document}

